\chapter{Traceability Matrix}\label{app:traceability}

This appendix traces the thesis's principal claims from motivation to evidence. For traceability, this
matrix uses four evaluation questions. RQ1 asks whether each modelled key substitution is detected and,
for an established conversation affected by an A1 or A5 substitution, whether the detecting client can notify its peer in band. RQ2
asks whether adding that detection avoids enlarging an individual lookup proof, and what the measured
proof sizes and local component timings indicate about the component cost of detection. RQ3 asks
whether offline catch-up avoids one proof per missed epoch and how the resulting proof size depends on
the two endpoint log sizes.
RQ4 asks whether the client verifier can integrate into the Signal Desktop codebase.
\Cref{tab:trace-rq} links each question to the corresponding design claim, mechanism, tests,
benchmarks, and evidence. \Cref{tab:trace-contrib} organizes the same evidence by the four
contributions of \cref{sec:intro-contributions}.

\begin{landscape}
\begin{table}[htbp]
\centering
\caption[Research-question traceability.]{Research question $\to$ claim $\to$ mechanism $\to$ test $\to$ benchmark $\to$ evidence.}
\label{tab:trace-rq}
\setlength{\extrarowheight}{0pt}
\renewcommand{\arraystretch}{0.9}
\setstretch{1}
\footnotesize
\begin{tabular}{@{}L{1.6cm}L{4.3cm}L{3.7cm}L{5.4cm}L{1.7cm}L{5.4cm}@{}}
\toprule
RQ & Design claim & Mechanism & Test or scenario & Benchmark & Evidence \\
\midrule
RQ1 & each modelled substitution is detected.\newline After honest establishment, detection of A1
or A5 also produces an in-band peer warning.\newline A3 is detected locally by the global
VRF-bound identity check when the directory is honest &
A1: dual-index warning scheme (\cref{sec:design-dual}).\newline A2: compound warning-censorship
signal and post-commit controls.\newline A3: global VRF-bound identity check, honest
directory.\newline A5: identity re-fetch and predecessor cross-signature. &
A1, A2, A3, and A5 scenarios (\cref{tab:tests-adversarial}) and dual-index derivation tests
(\cref{tab:tests-rust-prim}) &
M17 (clean workflows) & \cref{sec:eval-security}: the specified adversarial cases produce their expected outcomes.
M17 asserts a clean verdict without a warning for every listed clean workflow \\
RQ2 & detection does not enlarge an individual lookup proof. Across the tested large-directory
sizes, median full verification changes from \SI{133.6}{\micro\second} at $N=10^6$ to
\SI{139.2}{\micro\second} at $N=10^7$ &
combined-tree verification chain (\cref{ssec:design-verify}) &
prefix-tree, VRF, and Merkle-log tests (\cref{tab:tests-rust-prim,tab:tests-rust-server}) &
M1--M3, M7, M11, M16 & \cref{sec:eval-micro,ssec:eval-scale}: median proof sizes are
\num{817}\,B at $N=10^6$ and \num{945}\,B at $N=10^7$. M16 measures
\SI{27.8}{\micro\second} and \SI{31.1}{\micro\second} for reconstructed mismatch and clean
paths in the \texttt{duet} crate. The Signal Desktop path is not timed
(\cref{ssec:eval-detection}) \\
RQ3 & offline catch-up uses one consistency proof whose size is logarithmic in the endpoint log
sizes &
windowed consistency proof (\cref{sec:design-windowed,lem:windowed-size}) &
Merkle-log exhaustive-consistency and windowed-consistency tests
(\cref{tab:tests-rust-prim,tab:tests-rust-server}) &
M4 & \cref{ssec:eval-catchup}: at $W=720$, the one-proof construction contains \num{320}\,B of
raw consistency-path hashes, compared with \num{222528}\,B across the synthetic $W$-proof
baseline \\
RQ4 & the verifier integrates into Signal Desktop as a prototype warning path that does not block
session establishment &
Signal Desktop verifier and lifecycle hooks (\cref{sec:impl-signal}) &
client unit and wire-parity suites (\cref{sec:tests-ts}) and the live demonstration
(\cref{ssec:impl-demo}) &
Not applicable & \cref{sec:impl-signal,ssec:impl-demo}: the verifier and warning interface run in
the Signal Desktop fork, and the two-client demonstration exercises the integrated path \\
\midrule
Supporting & split-view equivocation is detected by comparing separately obtained signed heads &
independent-vantage signed-head comparison (\cref{alg:auditor}) &
the A4 scenario and auditor fork tests (\cref{sec:tests-rust}) &
M12 & \cref{ssec:eval-auditor}: M12 measures a local comparison after one head is stored. It does
not measure acquisition from an independent vantage \\
Supporting & continuous mode requires more directory writes than session mode under the assumed
session interval &
CEA continuous mode (\cref{ssec:design-modes}) &
M13 mode experiment (\cref{sec:bench-modes}) &
M13 & \cref{ssec:eval-modes}: the analytic ratio ranges from $24\times$ to $720\times$ across the
specified epoch counts and assumed session interval. These counts are not measured Signal
traffic \\
Supporting & a committed warning cannot be removed silently &
write-once policy, transition audit, and persistence evidence
(\cref{ssec:design-detection}) &
write-once and undeclared-removal tests (\cref{tab:tests-rust-server}) and transition audit and
deletion-evidence tests (\cref{tab:tests-rust-system}) &
None & \cref{sec:eval-security}: post-commit warning removal is detected by transition replay,
with persistence evidence providing an additional deletion record \\
\bottomrule
\end{tabular}
\end{table}
\end{landscape}

\begin{table}[ht]
\centering
\caption{Contribution kind $\to$ output $\to$ primary evidence.}
\label{tab:trace-contrib}
\setlength{\extrarowheight}{2pt}
\footnotesize
\begin{tabularx}{\textwidth}{@{}>{\raggedright\arraybackslash}p{2.1cm}>{\raggedright\arraybackslash}X>{\raggedright\arraybackslash}X@{}}
\toprule
Contribution & Output & Primary evidence \\
\midrule
Conceptual & dual-index warning scheme and ACKA-log instantiation (\cref{sec:design-dual,sec:design-acka}) & expected in-band verdicts in the listed dual-index warning cases. The implemented directory and auditor use ACKA's secure-log interface (\cref{sec:eval-security}) \\
Cryptographic & windowed consistency proof: RFC~9162 distant-head proof applied to offline monitoring, size bound (\cref{lem:windowed-size}) and the dual-index same-head consistency lemma as supporting analysis (\cref{lem:dualindex-consistency}) & same-head label consistency shown from proof soundness and collision resistance. \num{320}\,B raw path bytes at $W{=}720$, one proof for $W$ missed epochs, size logarithmic in endpoint log sizes (\cref{ssec:eval-catchup}) \\
Systems & the \texttt{duet} crate containing the directory server, cryptographic core, auditor
process, and benchmark harness, together with the Signal Desktop client component
(\cref{ch:impl}) & fixed vectors and live integration tests establish agreement between the crate
and client component, while the two-client demonstration exercises the integrated warning path
(\cref{ssec:impl-demo}) \\
Empirical & benchmark evaluation M1--M17 and the STRIDE/MITRE evaluation
(\cref{ch:eval,app:security}) & detection reuses the existing lookup-proof format, while median full
verification remains below \SI{150}{\micro\second} across the tested sizes
(\cref{sec:eval-micro}). The adversarial suite produces the expected outcome for every specified
detection scenario, using global VRF-bound identity comparison, the shared warning path, identity
re-fetch, predecessor cross-signature verification, and supplied signed-head comparison
(\cref{sec:eval-security}) \\
\bottomrule
\end{tabularx}
\end{table}

